\ECchapter{03}{Wind Energy}{How do engineers harvest the wind?}{ECGreen}

\ECObjectives{
\begin{itemize}[leftmargin=*,itemsep=1pt,topsep=1pt]
\item Explain how aerodynamic lift creates rotor torque.
\item Identify the main subsystems of a modern wind turbine.
\item Read and interpret a wind-turbine power curve.
\item Compare onshore, fixed-bottom offshore and floating offshore solutions.
\end{itemize}}

\begin{multicols}{2}

\begin{engineeringnews}
Modern offshore wind turbines can exceed 15~MW. Floating foundations open access to
 deeper-water sites, where winds are often stronger and steadier. They also create new
engineering challenges involving moorings, dynamic cables, maintenance and grid connection.
\end{engineeringnews}

\begin{ecarticle}
A wind turbine converts part of the kinetic energy of moving air into electrical energy.
Its blades are shaped like aerofoils. As air flows around a blade, the pressure on one side
becomes lower than on the other. This pressure difference creates a lift force. Because the
blade is inclined and rotating, part of this lift acts tangentially and produces a torque on
the rotor.

The power available in the wind is
\[
P_{\mathrm{wind}}=\frac{1}{2}\rho A v^3,
\]
where $\rho$ is air density, $A$ the swept area and $v$ wind speed. The cubic dependence on
$v$ is essential: doubling wind speed multiplies the available power by eight. This is why
careful site measurement and long-term wind statistics are more important than a single
instantaneous reading.

The rotor drives either a gearbox and a high-speed generator or a larger direct-drive
generator. Power electronics adapt the variable electrical output to the frequency and
voltage of the grid. Pitch control changes the angle of the blades to regulate power and
limit loads. Yaw control turns the nacelle so that the rotor faces the wind.

A wind turbine does not produce its rated power all the time. Below the cut-in speed, it
remains stopped. Between cut-in and rated speed, output rises rapidly. Above rated speed,
control systems limit the power. At the cut-out speed, the turbine shuts down to avoid
excessive mechanical loads.
\end{ecarticle}

\begin{engineeringfocus}
\textbf{Main subsystems}
\begin{enumerate}[leftmargin=*,itemsep=1pt]
\item rotor, blades and hub;
\item main shaft and bearings;
\item gearbox or direct-drive generator;
\item converter and transformer;
\item pitch and yaw actuators;
\item tower, foundations and monitoring system.
\end{enumerate}

A reliable design must combine high energy capture with acceptable fatigue loads, safe
maintenance and a lifetime of several decades.
\end{engineeringfocus}

\begin{tcolorbox}[title=\sffamily\bfseries Did You Know?,colback=yellow!10,colframe=ECOrange]
No turbine can capture all the power in the wind. The theoretical upper limit is the
\textbf{Betz limit}: 59.3\% of the kinetic power crossing the rotor area.
\end{tcolorbox}

\begin{tcolorbox}[title=\sffamily\bfseries Key Figures,colback=ECLightBlue,colframe=ECBlue]
\begin{tabularx}{\linewidth}{@{}Xr@{}}
Rotor diameter & 120--250 m\\
Rated power & 3--18 MW\\
Cut-in speed & 3--4 m/s\\
Rated speed & 11--15 m/s\\
Cut-out speed & 20--25 m/s\\
Typical capacity factor & 30--60\%
\end{tabularx}
\end{tcolorbox}

\columnbreak

\begin{tcolorbox}[title=\sffamily\bfseries Wind-turbine architecture,colback=white,colframe=ECBlue]
\centering
\begin{tikzpicture}[scale=.68,transform shape,>=Latex,font=\scriptsize]
\draw[very thick,ECGreen] (0,0) -- (0,3.4);
\draw[fill=ECGray,draw=ECBlue,thick] (-.7,3.4) rectangle (1.8,4.1);
\draw[fill=ECBlue!10,draw=ECBlue,thick] (1.8,3.52) rectangle (2.8,3.98);
\draw[thick] (-.7,3.75)--(-1.45,3.75);
\filldraw[ECOrange] (-1.45,3.75) circle (.12);
\foreach \a in {0,120,240}{
  \draw[very thick,ECGreen,rotate around={\a:(-1.45,3.75)}]
  (-1.45,3.75)--(-1.45,5.0);
}
\draw[->,thick] (2.9,3.75)--(4.0,3.75) node[right]{electricity};
\node[align=center] at (.55,3.75){gearbox /\\generator};
\node[align=center] at (0,-.35){tower and\\foundation};
\node[left] at (-1.5,4.7){blades};
\node[above] at (.55,4.1){nacelle};
\end{tikzpicture}
\end{tcolorbox}

\begin{tcolorbox}[title=\sffamily\bfseries Typical power curve,colback=white,colframe=ECGreen]
\centering
\begin{tikzpicture}
\begin{axis}[
width=\linewidth,height=47mm,
xlabel={Wind speed (m/s)},ylabel={Power (MW)},
xmin=0,xmax=26,ymin=0,ymax=3.3,
grid=major,
tick label style={font=\scriptsize},
label style={font=\scriptsize},
]
\addplot[very thick,ECBlue,smooth] table[x=wind,y=power,col sep=comma]
{data/EC03/wind_power_curve.csv};
\addplot[dashed,ECOrange] coordinates {(3,0) (3,3.2)};
\addplot[dashed,ECOrange] coordinates {(15,0) (15,3.2)};
\addplot[dashed,ECOrange] coordinates {(25,0) (25,3.2)};
\end{axis}
\end{tikzpicture}

{\scriptsize The cut-in, rated and cut-out regions determine how the turbine operates safely.}
\end{tcolorbox}

\begin{tcolorbox}[title=\sffamily\bfseries Engineering Insight,colback=white,colframe=ECOrange]
A larger rotor does more than increase rated power. It also captures more energy at moderate
wind speeds, which can improve annual production. However, longer blades experience higher
loads and are more difficult to manufacture, transport and recycle.
\end{tcolorbox}

\begin{vocabulary}
\begin{tabularx}{\linewidth}{@{}lX@{}}
blade & pale\\
hub & moyeu\\
gearbox & multiplicateur\\
yaw & orientation de la nacelle\\
pitch & calage des pales\\
swept area & surface balayée\\
capacity factor & facteur de charge
\end{tabularx}
\end{vocabulary}

\begin{questions}
\item Why does wind power depend on the cube of wind speed?
\item How does aerodynamic lift create rotor torque?
\item What is the purpose of pitch control?
\item What is the purpose of yaw control?
\item Interpret the main regions of the power curve.
\item Why can offshore wind achieve a higher capacity factor than onshore wind?
\end{questions}

\begin{challenge}
A coastal region plans a 500~MW wind farm. Compare an onshore solution, a fixed-bottom
offshore solution and a floating offshore solution. Consider energy yield, foundations,
maintenance, grid connection, environmental impact and public acceptance. Use a weighted
decision table and present a justified recommendation.
\end{challenge}

\begin{applicationexercise}
A wind turbine has a rotor diameter of 100~m and operates in air of density
$1.225~\mathrm{kg\,m^{-3}}$ at a wind speed of 10~m/s. Using
$P=\tfrac12\rho A v^3 C_p$ with $C_p=0.44$, estimate the aerodynamic power captured by the
rotor. Compare $C_p$ with the Betz limit, then estimate the electrical power delivered to the
grid if the mechanical and electrical efficiency is 90\%.
\end{applicationexercise}

\begin{speaking}
Debate the statement: ``Large wind farms should be built mainly offshore.'' Use technical,
economic and environmental arguments, and respond to at least one opposing argument.
\end{speaking}

\end{multicols}

\subsection*{Sources and further reading}
\ECsource{Global Wind Energy Council}{https://gwec.net}\quad
\ECsource{IEA -- Wind}{https://www.iea.org/energy-system/renewables/wind}
